% Readout SAE transfer and model-family comparison appendix.

\section{Readout SAE Transfer and Model-Family Comparisons}
\label{app:readout-factorizer-transfer-comparisons}

This appendix reports how the selected readout SAE recipe transfers across the
smaller Qwen/Gemma suite and targeted cross-family or post-training rows.

\subsection{Transfer to Smaller Models}
\label{app:readout-factorizer-small-model-transfer}

The fixed Qwen3.5-9B recipe was then transferred to four smaller models. The
transfer tests both whether a fixed active budget reconstructs
lower-dimensional unembedding rows more effectively, and whether the higher
$k=256$ fidelity regime remains useful away from 9B. As above, native $k=128$
and native $k=256$ rows are reported as separate operating regimes; truncation
comparisons are handled as capacity controls.

\paragraph{Coverage.}
The 5k screen covered all 24 model/width/$k$ cells
(Fig.~\ref{fig:app-k-small-model-pilot-grid}). In every model, the largest
screening rowEV occurred at 32$\times$, $k=256$: Qwen3.5-0.8B $0.812$,
Qwen3.5-2B $0.765$, Gemma-4-E2B $0.711$, and Gemma-4-E4B $0.742$. These screening
values select continuations; the 20k rows in
Table~\ref{tab:app-k-small-model-converged-ledger} supply the reported
dictionary settings.

\begin{figure}[!tbp]
    \centering
    \ReadoutSAEGraphic[height=0.26\textheight]{appendix_k_small_model_pilot_grid.pdf}
    \caption{
        Small-model screening grid. Each model was evaluated at 8$\times$,
        16$\times$, and 32$\times$ under native $k=128$ and $k=256$; final
        claims use the native 20k rows in
        Table~\ref{tab:app-k-small-model-converged-ledger}.
    }
    \label{fig:app-k-small-model-pilot-grid}
\end{figure}

\begin{table*}[!tbp]
\caption{Twenty-thousand-step small-model results. The 16$\times$, $k=128$
rows instantiate the strict-budget regime; the 32$\times$, $k=256$ rows
instantiate the high-fidelity regime. Gemma top1/KL are
softcap-correct.}
\label{tab:app-k-small-model-converged-ledger}
\centering
\footnotesize
\setlength{\tabcolsep}{4pt}
\renewcommand{\arraystretch}{1.08}
\begin{tabularx}{\textwidth}{@{}p{0.14\textwidth}p{0.12\textwidth}rrrrY@{}}
\toprule
Model & Dictionary setting & rowEV & top1 & KL & dead/rare & Role \\
\midrule
Qwen3.5-0.8B
& 16$\times$, $k=128$
& $0.760$ & $0.844$ & $0.277$ & $0.000/0.001$
& strict-budget transfer result \\
\addlinespace
Qwen3.5-0.8B
& 32$\times$, $k=256$
& $\mathbf{0.877}$ & $0.891$ & $0.135$ & $0.000/0.001$
& high-fidelity transfer result \\
\addlinespace
Qwen3.5-2B
& 16$\times$, $k=128$
& $0.712$ & $0.858$ & $0.261$ & $0.000/0.004$
& strict-budget transfer \\
\addlinespace
Qwen3.5-2B
& 32$\times$, $k=256$
& $0.847$ & $0.887$ & $0.136$ & $0.000/0.010$
& high-fidelity transfer \\
\addlinespace
Gemma-4-E2B
& 16$\times$, $k=128$
& $0.714$ & $0.623$ & $1.94$ & $0.001/0.031$
& softcap-correct fidelity \\
\addlinespace
Gemma-4-E2B
& 32$\times$, $k=256$
& $0.834$ & $0.333$ & $6.37$ & $0.000/0.038$
& high rowEV; lower replacement fidelity \\
\addlinespace
Gemma-4-E4B
& 16$\times$, $k=128$
& $0.693$ & $0.669$ & $1.82$ & $0.012/0.040$
& softcap-correct fidelity \\
\addlinespace
Gemma-4-E4B
& 32$\times$, $k=256$
& $0.827$ & $0.736$ & $1.22$ & $0.002/0.078$
& fidelity row; softcap-correct \\
\bottomrule
\end{tabularx}
\end{table*}

\begin{table}[!tbp]
\caption{Per-$k$ evaluation curve for the completed Gemma-4-E4B 32$\times$,
$k=256$ convergence.}
\label{tab:app-k-gemma-e4b-k-curve}
\centering
\footnotesize
\setlength{\tabcolsep}{6pt}
\renewcommand{\arraystretch}{1.10}
\begin{tabular}{@{}lrrr@{}}
\toprule
$k_{\mathrm{eval}}$ & rowEV & top1 & KL (bits) \\
\midrule
$64$ & $0.113$ & $0.008$ & $38.5$ \\
$128$ & $0.568$ & $0.037$ & $31.6$ \\
$192$ & $0.774$ & $0.321$ & $10.8$ \\
$256$ native & $0.827$ & $0.736$ & $1.22$ \\
$384$ & $0.684$ & $0.720$ & $1.37$ \\
\bottomrule
\end{tabular}
\end{table}

The native $k=256$ evaluation is the reported dictionary setting; the $k=384$ row
is an over-budget comparison. KL is softcap-correct, while rowEV and top1 are
cap-independent / argmax-invariant.

\begin{figure}[!tbp]
    \centering
    \ReadoutSAEGraphic[height=0.28\textheight]{appendix_k_cross_model_scaling.pdf}
    \caption{
        Qwen scaling summary. Panel A reports the Qwen rows used for the
        scaling claim; Panel B reports rare feature rate as a usage statistic.
        Additional model-family rows are reported in the surrounding tables.
    }
    \label{fig:app-k-cross-model-scaling}
\end{figure}

\ReadoutSAEFloatBarrier

Across the native Qwen dictionary settings in
Fig.~\ref{fig:app-k-cross-model-scaling}, additional capacity generally improves
both rowEV and top-1 agreement. The selected strict-budget transfer point is
Qwen3.5-0.8B 16$\times$, $k=128$; the selected high-fidelity
small-model point is Qwen3.5-0.8B 32$\times$, $k=256$, with Qwen3.5-2B giving
the same qualitative pattern. The scaling figure focuses on Qwen so that model
size is varied within a single family; the surrounding tables evaluate the same
recipe under additional readout geometries. Rare feature rate is reported as a
usage statistic: it stays low for
the smaller Qwen 32$\times$, $k=256$ rows and rises to $0.292$ for the 9B
32$\times$, $k=256$ capacity point.
Qwen3.5-2B has slightly higher strict-budget top1, so the 0.8B selection is a
joint reconstruction-and-fidelity dictionary setting with top1 as one component.

\begin{table*}[!tbp]
\caption{Qwen3.5-0.8B 32$\times$, $k=256$ seed-variation window. Seed 0 is the
archived dictionary setting; seeds 1 and 2 repeat the same recipe.}
\label{tab:app-k-qwen08b-seed-window}
\centering
\scriptsize
\setlength{\tabcolsep}{4pt}
\renewcommand{\arraystretch}{1.08}
\begin{tabular}{@{}lrrrrl@{}}
\toprule
Seed & rowEV & top1 & KL & rare & Wall \\
\midrule
0 (archived) & $0.877$ & $0.891$ & $0.135$ & $0.001$ & -- \\
1 & $0.886$ & $0.875$ & $0.104$ & $0.134$ & 39 min \\
2 & $0.887$ & $0.879$ & $0.105$ & $0.142$ & 51 min \\
\midrule
mean $\pm$ std & $0.883 \pm 0.006$ & $0.882 \pm 0.009$ &
$0.115 \pm 0.018$ & $0.092 \pm 0.080$ & -- \\
\bottomrule
\end{tabular}
\end{table*}

The seed window shows that the headline replacement metrics are stable under
this recipe: rowEV, top1, and KL have standard deviation at most $0.02$. Rare
feature rate varies more across seeds, so we report it as a usage statistic and
use the seed-window mean when making seed-robust statements about that tail.

\ReadoutSAEFloatBarrier
\subsection{Cross-Family 8B Comparison (Ministral)}
\label{app:readout-factorizer-ministral-comparison}

Ministral-3-8B-Base is included as a cross-family 8B-scale comparison. It has the
same hidden width $d_{\mathrm{model}}=4096$ as Qwen3.5-9B but an untied
\texttt{lm\_head} and no post-readout softcap (readout-extraction
reconstruction relative error $\approx 3{\times}10^{-6}$, an order of magnitude
cleaner than the bf16 Qwen floor). The same transferred TopK recipe was used
at the matched widths.

\begin{table*}[!tbp]
\caption{Cross-family and post-training metrics under the
transferred TopK recipe. ``Strict'' denotes native $k=128$; ``Fidelity''
denotes native $k=256$.}
\label{tab:app-k-cross-family-stress-converged}
\centering
\scriptsize
\setlength{\tabcolsep}{4pt}
\renewcommand{\arraystretch}{1.00}
\begin{tabularx}{0.96\textwidth}{@{}p{0.18\textwidth}p{0.15\textwidth}p{0.12\textwidth}rrrrrY@{}}
\toprule
Model & Dictionary setting & Width / $k$ & rowEV & top1 & KL & dead & rare & Role \\
\midrule
\shortstack[l]{Ministral\\3-8B}
& Strict
& 16$\times$, $k=128$
& $0.806$ & $0.885$ & $0.130$ & $0.026$ & $0.456$
& strict-budget comparison \\
\addlinespace
\shortstack[l]{Ministral\\3-8B}
& Fidelity
& 32$\times$, $k=256$
& $0.888$ & $0.904$ & $0.087$ & $0.018$ & $0.529$
& fidelity comparison \\
\addlinespace
\shortstack[l]{R1-Distill\\Qwen-7B}
& Strict
& 16$\times$, $k=128$
& $0.709$ & $0.695$ & $0.777$ & $0.016$ & $0.298$
& strict-budget comparison \\
\addlinespace
\shortstack[l]{R1-Distill\\Qwen-7B}
& Fidelity
& 32$\times$, $k=256$
& $0.844$ & $0.760$ & $0.489$ & $0.016$ & $0.378$
& fidelity gap \\
\addlinespace
\shortstack[l]{R1-Distill\\Llama-8B}
& Strict
& 16$\times$, $k=128$
& $0.796$ & $0.725$ & $0.536$ & $0.009$ & $0.356$
& strict-budget comparison \\
\addlinespace
\shortstack[l]{R1-Distill\\Llama-8B}
& Fidelity
& 32$\times$, $k=256$
& $0.888$ & $0.754$ & $0.434$ & $0.011$ & $0.392$
& fidelity gap (rowEV-matched, top1/KL apart) \\
\bottomrule
\end{tabularx}
\end{table*}

Both Ministral dictionary settings have healthy usage and decision-fidelity
metrics (dead $< 0.03$, KL $< 0.14$ bits).
Decision fidelity (top1, KL) falls in the high-fidelity range of the matched
32$\times$, $k=256$ comparison rows, supporting the claim that the
row-factorization recipe transfers beyond the Qwen family when the readout is
geometrically clean (untied, softcap-free). The rare feature rate
is high at both points; we report it as a usage statistic, following the same
convention as the Qwen3.5-9B
32$\times$, $k=256$ row.

\ReadoutSAEFloatBarrier
\subsection{R1-Distilled Comparison Rows}
\label{app:readout-factorizer-r1-distill-comparison}

R1-Distill-Qwen-7B is included as a post-training comparison row for comparing
rowEV and decision fidelity. It is a Qwen-derived 7B readout with an untied
\texttt{lm\_head} and no post-readout softcap, giving a clean comparison under the same
transferred TopK recipe family at the strict-budget and high-fidelity dictionary
settings in
Table~\ref{tab:app-k-cross-family-stress-converged}.

The 32$\times$, $k=256$ R1 row preserves high row reconstruction, within
$0.044$ rowEV of the matched Ministral point, while top1 and KL separate from
the base-model comparison. This gives a third pattern on the rowEV vs.\ fidelity
axis: row reconstruction and decision-level readout replacement are distinct
evaluation quantities.

R1-Distill-Llama-8B adds a second reasoning-distilled checkpoint from a
non-Qwen base family at both dictionary settings. At the strict budget, its
rowEV is close to the Ministral strict-budget row ($0.796$ vs.\ $0.806$), with
different top1/KL behavior ($0.725/0.536$ vs.\ $0.885/0.130$). At the fidelity
setting the row-reconstruction matches the Ministral fidelity row exactly
($0.888$ vs.\ $0.888$), while top1/KL remain separated ($0.754/0.434$ vs.\
$0.904/0.087$). The pattern is consistent across both base families
(Qwen2.5-Math-derived \citep{yang2024qwen25math} and Llama-3.1-derived
\citep{grattafiori2024llama3herd}): reasoning-distilled readouts can preserve
row reconstruction while requiring separate decision-fidelity
measurement. Appendix~\ref{app:readout-norm-tails} reports the row-norm tail
analysis that helps explain this separation.

\paragraph{Attribution of the post-training effect.}
The R1-distilled checkpoints differ from the base-model rows in architecture,
tokenizer, and pretraining lineage as well as post-training. The repeated pattern
across two distillation lineages makes post-training a plausible contributor; a
matched base-vs.-distill ablation within one architecture would further isolate
the effect.

\paragraph{R1 contrastive readout-target fidelity.}
We also ran the held-out contrastive readout-target fidelity check on the two
R1-distilled checkpoints. Table~\ref{tab:app-k-r1-qfid} reports these results
using the same bank version as the five-model run. Sign agreement remains close
to the five-model baseline; the fraction below the residual threshold is
comparable in the fidelity setting and lower under the stricter setting.

\begin{table}[!tbp]
\caption{Held-out contrastive readout-target fidelity metrics for the R1-distilled
checkpoints under the same contrast-bank version as the five-model run.
``sign'' is contrastive sign agreement; ``median resid.'' is the median relative
residual $|\Delta\text{margin}|/|\text{exact margin}|$; ``resid.\(<0.5\)''
is the fraction of held-out readout targets with relative residual below $0.5$.}
\label{tab:app-k-r1-qfid}
\centering
\scriptsize
\setlength{\tabcolsep}{2.5pt}
\renewcommand{\arraystretch}{1.08}
\begin{tabularx}{\columnwidth}{@{}>{\RaggedRight\arraybackslash}X>{\RaggedRight\arraybackslash}p{0.28\columnwidth}ccc@{}}
\toprule
Model & \makecell[l]{Setting\\(cfg)} & sign & \makecell{median\\resid.} & \makecell{resid.\\$<0.5$} \\
\midrule
\shortstack[l]{Ministral\\3-8B-Base}  & \makecell[l]{fidelity\\(32$\times$, $k=256$)} & $0.949$ & $0.079$ & $0.831$ \\
\shortstack[l]{Ministral\\3-8B-Base}  & \makecell[l]{strict\\(16$\times$, $k=128$)}   & $0.947$ & $0.102$ & $0.812$ \\
\addlinespace
\shortstack[l]{R1-Distill\\Llama-8B} & \makecell[l]{fidelity\\(32$\times$, $k=256$)} & $0.926$ & $0.146$ & $0.782$ \\
\shortstack[l]{R1-Distill\\Llama-8B} & \makecell[l]{strict\\(16$\times$, $k=128$)}   & $0.907$ & $0.193$ & $0.752$ \\
\addlinespace
\shortstack[l]{R1-Distill\\Qwen-7B}  & \makecell[l]{fidelity\\(32$\times$, $k=256$)} & $0.921$ & $0.163$ & $0.738$ \\
\shortstack[l]{R1-Distill\\Qwen-7B}  & \makecell[l]{strict\\(16$\times$, $k=128$)}   & $0.890$ & $0.233$ & $0.700$ \\
\midrule
\multicolumn{2}{l}{5-model baseline}
                                       & $0.934$ & --      & $0.773$ \\
\bottomrule
\end{tabularx}
\end{table}

\ReadoutSAEFloatBarrier
\subsection{Readout Row-Norm Tail Analysis}
\label{app:readout-norm-tails}
\label{sec:app-readout-norm-tails}

The TopK recipe operates on row-centered, row-normalized $\WU$ rows, which
treats all rows uniformly in $L_2$. When the underlying row-norm distribution
is heavy-tailed, however, a small per-row reconstruction error is amplified
disproportionately on the high-norm rows — and those are precisely the rows
that disproportionately affect top1 / KL after the softmax. We therefore
report row-norm summary statistics for the base-model and reasoning-distilled
checkpoints as a usage analysis for the separation between rowEV and fidelity
observed in Appendix~\ref{app:readout-factorizer-r1-distill-comparison}.

\begin{table}[!tbp]
\caption{$\WU$ row-norm statistics for the base-model and reasoning-distilled
checkpoints, computed on the readout extraction artifacts used to train each
SAE. ``CoV'' is the coefficient of variation $\sigma/\mu$. ``max/median''
summarises the heavy-tail end.}
\label{tab:app-k-row-norm-tail}
\centering
\scriptsize
\setlength{\tabcolsep}{2.5pt}
\renewcommand{\arraystretch}{1.08}
\resizebox{\columnwidth}{!}{%
\begin{tabular}{@{}lrrrrrr@{}}
\toprule
Model & min & p50 & max & mean & CoV & max/med \\
\midrule
Qwen3.5-9B (base)         & $0.379$ & $0.980$ & $1.545$ & $0.983$ & $0.120$ & $1.58$ \\
Ministral-3-8B (base)     & $0.250$ & $0.447$ & $0.684$ & $0.449$ & $0.116$ & $1.53$ \\
\addlinespace
R1-Distill-Qwen-7B (sft)  & $0.573$ & $1.054$ & $1.756$ & $1.029$ & $0.159$ & $1.67$ \\
R1-Distill-Llama-8B (sft) & $0.424$ & $0.897$ & $1.787$ & $0.904$ & $0.139$ & $1.99$ \\
\bottomrule
\end{tabular}
}
\end{table}

Two effects are visible. First, the upper tail is heavier on the distilled
checkpoints: max/median is $1.5$--$1.6$ for the base-model rows and
$1.7$--$2.0$ for the distilled rows, with R1-Distill-Llama-8B the most
extreme. Second, the lower tail is compressed: minimum row norms rise
from $0.25$--$0.38$ on the base-model rows to $0.42$--$0.57$ on the distilled
rows. In these comparison rows, the row-norm distribution is wider for the
reasoning-distilled checkpoints than for the base-model rows. The coefficient of
variation tracks this ($0.12$ on base models vs $0.14$--$0.16$ on distilled
models).

The row-centered-and-normalized SAE recipe therefore operates on a heavier-tailed
norm distribution for the distilled rows, which is consistent with the
observed pattern: row reconstruction (which is uniform in $L_2$) is preserved,
while decision fidelity (which is exponential in the logit error and dominated
by the heaviest rows) separates from rowEV. This characterizes a second-order
property of these reasoning-distilled comparison rows and motivates norm-aware
SAE recipes as a natural extension for those targets.

\ReadoutSAEFloatBarrier
